\section{Related literature}
The present work lies at the intersection of stochastic transportation planning, green logistics under carbon regulation, fractional calculus with long memory, and computational methods for risk-averse stochastic optimization. The purpose of this section is not only to position the model, but also to clarify which nearby strands of work provide useful algorithmic or empirical benchmarks.

\subsection{Stochastic transportation and recourse planning}
The classical transportation problem has been generalized to multi-period, capacity-constrained, and uncertainty-aware settings in which here-and-now decisions are separated from scenario-dependent recourse. Foundational references in stochastic programming explain how scenario trees, recourse functions, and sample-based approximations support such models \cite{birge2011,shapiro2009}. The present paper follows this tradition, but departs from standard formulations by replacing the usual first-order state recursion with a fractional balance law. As a result, recourse decisions affect not only the current period but also the weighted memory of past imbalance.

\subsection{Green logistics, carbon regulation, and allowance trading}
Environmental performance is now central to transportation modeling. The pollution-routing literature and later reviews of green vehicle-routing show that emissions considerations reshape routing, mode choice, and network design \cite{bektas2011,asghari2021}. In supply-chain design, robust and stochastic models have been used to study carbon policies, including cap-and-trade structures \cite{homayouni2023carbon,qu2021}. A particularly relevant recent comparison is the multi-stage green supply-chain design model of Mirzaee et al., which studies disruptions, resilience, and cap-and-trade under two-stage stochastic optimization \cite{mirzaee2024green}. Relative to that line of work, the present paper differs in two ways: it embeds memory directly in the operational state dynamics, and it models allowance trading as an explicit recourse channel rather than relying only on aggregate cap or penalty effects.

A second nearby contribution is the CVaR-based two-stage intermodal transport model of Camur et al., which studies uncertain demand and spot train capacities in a green road-rail setting \cite{gbadegoye2025intermodal}. That paper is important as an empirical logistics comparator because it combines transportation cost, emissions, and CVaR in a modern stochastic design. The key distinction is again structural: the present model introduces Caputo memory and endogenous allowance trading, whereas the intermodal model remains integer-order and focuses on intermodal capacity uncertainty.

\subsection{Fractional calculus and long-memory operational systems}
Fractional calculus is a standard tool for systems with hereditary effects, anomalous diffusion, and persistent dissipation. Foundational references include Podlubny, Caputo, and Diethelm \cite{podlubny1999,caputo1967,diethelm2010}. In operations and logistics, however, direct use of fractional dynamics remains limited. Existing work has used fractional formulations in inventory and adaptive supply-chain control, indicating that long-memory effects can matter when deterioration, disruption recovery, or congestion persistence are significant \cite{cuong2023,haque2026}. The current paper extends this logic to a transportation model with uncertain demand, service protection, and carbon-market recourse.

\subsection{Discretization, fast memory approximation, and computational tractability}
A central computational challenge in fractional optimization is the dense temporal coupling induced by the Caputo derivative. The L1 discretization remains attractive because of its transparency and direct compatibility with finite-horizon optimization models \cite{lin2007}. However, L1 is only one point in a broader numerical-design space. Fang et al. develop a fast higher-order approximation based on the $L2$-$1_{\sigma}$ formula combined with exponential-sum approximation, reducing both memory usage and computational cost for Caputo-type operators over longer horizons \cite{zhang2022fast}. This literature is directly relevant here: if the transportation horizon or scenario count grows, fast convolution and compressed-memory approximations become natural candidates for improving scalability beyond the current benchmark.

\subsection{Risk aversion, multistage models, and online CVaR methods}
Conditional value-at-risk is widely used in stochastic optimization because it captures tail severity while remaining tractable in continuous settings \cite{rockafellar2000}. Beyond the classical CVaR linearization, recent work shows that risk-averse transportation and network models can be framed in both two-stage and multistage ways. Yu and Shen study risk-averse two-stage and multistage stochastic facility-location models using expected conditional risk measures and derive performance gaps and computational enhancements between static and adaptive formulations \cite{yu2025multistage}. Their results are useful context for the present paper because they show that the choice between fixed and adaptive structures is itself a substantive modeling decision under risk.

At the algorithmic level, Madavan et al. analyze online primal-dual methods for CVaR-constrained and CVaR-penalized stochastic optimization \cite{madavan2021cvar}. Their work is not a direct replacement for progressive hedging in the current mixed-integer setting, but it provides an important alternative perspective: large-sample risk-averse optimization can also be approached through stochastic first-order updates rather than scenario decomposition. This is relevant when one wants to move beyond medium-scale SAA toward streaming or data-rich settings.

\subsection{Scenario quality and sample-based approximation}
Because the present manuscript relies on sample-average approximation, the quality of the scenario sample matters. Classical references discuss SAA consistency and policy aggregation \cite{rockafellar1991,shapiro2009}. More recently, Milz et al. establish laws of large numbers for sample-based approximations of law-invariant risk functionals and show that randomized quasi-Monte Carlo methods can reduce bias and root mean square error relative to plain Monte Carlo in risk-averse settings \cite{melnikov2025rqmc}. This literature strengthens the justification for reporting out-of-sample performance, stability across sample sizes, and, in future work, potentially using RQMC scenario generation to improve estimator efficiency.

\subsection{Research gap and contribution positioning}
The literature therefore suggests a clear gap. Stochastic transportation and green logistics models capture uncertainty and environmental trade-offs, but typically remain memoryless. Fractional models capture persistent memory, but are more common in dynamical systems and inventory theory than in transport optimization. Existing carbon-aware logistics studies provide useful comparators, yet they generally lack an explicit fractional state and endogenous allowance trading. The contribution of the present paper is to integrate these elements into one tractable framework while making explicit where future computational improvements should come from: stronger experimental baselines, richer out-of-sample assessment, and faster memory approximations for longer planning horizons.
